% !TEX TS-program = pdflatex
% !TEX encoding = UTF-8 Unicode
\documentclass[11pt,a4paper]{article}

\usepackage[utf8]{inputenc}
\usepackage[T1]{fontenc}
\usepackage{lmodern}
\usepackage[top=0.8in, bottom=0.8in, left=1in, right=1in]{geometry}
\usepackage{hyperref}
\usepackage{xcolor}
\usepackage{setspace}

% Colors
\definecolor{accentblue}{HTML}{1A73E8}

% Remove page number
\pagestyle{empty}

% Hyperlink setup
\hypersetup{
  colorlinks=true,
  linkcolor=accentblue,
  urlcolor=accentblue
}

% Reduce paragraph spacing
\setlength{\parskip}{0.5em}
\setlength{\parindent}{0pt}

\begin{document}

\begin{center}
  {\Large\bfseries LETTER OF RECOMMENDATION}
\end{center}

\vspace{0.3em}

Hanoi, March 23, 2026

\textbf{To:} Board of Directors / Human Resources Department\\
\textbf{Company:} VinRobotics Corporation (VinMotion)

\noindent\rule{\textwidth}{0.4pt}

Dear Sir/Madam,

I am \textbf{Assoc. Prof. Dr. Pham Duc An}, Vice Head of the Mechatronics Department and Head of the Automated Equipment Research Group at the School of Mechanical Engineering, Hanoi University of Science and Technology (HUST). I am writing this letter to wholeheartedly recommend \textbf{Le Duc Nguyen} --- a final-year student in Mechatronics Engineering (Class of 2021--2026) at our university --- for the \textbf{Engineering Internship} position at VinRobotics Corporation.

I have known Nguyen since he enrolled in several specialized courses under my instruction, including \textit{Autonomous Mobile Robots}, \textit{Mechatronic System Design}, and \textit{Image Processing}. I was also his direct supervisor for his graduation thesis, titled \textbf{``Multi-AGV Coordination System Using Reinforcement Learning''} --- a project that he completed with excellent results in my laboratory. Through teaching and working closely with him, I have a well-grounded basis to assess both his professional capabilities and personal qualities.

Nguyen is one of the most outstanding students I have supervised in recent years. In his graduation thesis, he demonstrated exceptional independent research ability by successfully developing a multi-AGV coordination system utilizing Reinforcement Learning, integrating motion control and simulation on the ROS platform. He showed a remarkably quick grasp of interdisciplinary knowledge --- from control theory and artificial intelligence to systems programming --- and applied them systematically to solve practical problems. This project is highly applicable and directly relevant to logistics and smart manufacturing --- domains that I understand VinRobotics is actively pursuing.

Beyond his thesis, Nguyen has actively participated in various research activities and engineering competitions. He was awarded \textbf{Third Prize at the DENSO Factory Hacks 2023} as a team leader, and received the \textbf{Scientific Research \& Innovation Award} from HUST --- reflecting his proactive spirit and practical problem-solving capabilities.

In terms of personal qualities, Nguyen is a \textbf{proactive, ambitious, and highly responsible} student. He is always willing to embrace new challenges and does not hesitate to self-study cutting-edge technologies. In the lab, he demonstrated strong independent work capability while being highly collaborative in team settings. He possesses clear systems thinking and the ability to articulate ideas coherently --- qualities essential in a research and product development environment.

With a solid foundation in \textbf{Robotics, Reinforcement Learning, ROS, and embedded systems}, combined with practical experience in \textbf{multi-robot coordination, industrial AI, and IoT}, I firmly believe that Nguyen is an excellent fit for VinRobotics. I \textbf{strongly recommend} Le Duc Nguyen to your esteemed organization and am happy to provide further information if needed.

\vspace{1em}

Sincerely,

\vspace{1.5em}

\textbf{Assoc. Prof. Dr. Pham Duc An}\\
Senior Lecturer --- Vice Head of Mechatronics Department\\
Head of Automated Equipment Research Group\\
School of Mechanical Engineering, Hanoi University of Science and Technology\\
Email: \href{mailto:an.phamduc@hust.edu.vn}{an.phamduc@hust.edu.vn} \quad|\quad Phone: (+84) 868 040 770

\end{document}
